% Compile from the FSI repository root:
%   pdflatex -output-directory docs/resources/benchmarks/abramson1965_report/build \
%     docs/resources/benchmarks/abramson1965_report/abramson1965_current_report.tex
% Or: docs/resources/benchmarks/abramson1965_report/build_report.sh
\documentclass[11pt,a4paper]{article}

\usepackage[utf8]{inputenc}
\usepackage[T1]{fontenc}
\usepackage[italian]{babel}
\usepackage{geometry}
\usepackage{graphicx}
\usepackage{subcaption}
\usepackage{booktabs}
\usepackage{siunitx}
\usepackage{amsmath}
\usepackage{amssymb}
\usepackage{float}
\usepackage{xcolor}
\usepackage{hyperref}
\usepackage{pgfplots}
\pgfplotsset{compat=1.18}
\pgfplotsset{every axis/.append style={grid=both, grid style={gray!20}, tick align=outside}}

\makeatletter
\newcommand{\fsiroot}{}
\newcommand{\fsipath}[1]{\fsiroot#1}
\IfFileExists{output_data/ABRAMSON1965/vlambda_data.csv}{%
  \def\fsiroot{}%
}{%
  \IfFileExists{../../../../output_data/ABRAMSON1965/vlambda_data.csv}{%
    \def\fsiroot{../../../../}%
  }{%
    \PackageWarning{abramson1965}{Cannot locate FSI repo root; figures/CSV may be missing.}%
  }%
}
\makeatother

\geometry{margin=24mm}
\sisetup{detect-all=true, per-mode=symbol}
\DeclareSIUnit{\knunit}{knots}
\DeclareSIUnit{\cps}{cps}
\DeclareSIUnit{\ftunit}{ft}

\title{Stato dell'analisi idro-aeroelastica dell'ala ABRAMSON1965}
\author{Report tecnico di lavoro}
\date{4 giugno 2026}

\begin{document}
\maketitle

\begin{abstract}
Questo report documenta la simulazione numerica del modello di idrofoil
ABRAMSON1965, confrontando gli input del caso di calcolo con i parametri del
lavoro originale di Abramson e Ransleben (1965), \emph{An Experimental Investigation of
Flutter of a Fully Submerged Subcavitating Hydrofoil}.
Il riferimento sperimentale e' la rottura per flutter a \SI{48.1}{\knunit} con
frequenza accoppiata di circa \SI{17.5}{\cps} e $(V/b\omega)_F=1{,}48$.
La linea di calcolo di riferimento e' \textbf{Roger RFA + DLM} con
\textbf{massa aggiunta da strip theory} sui primi due modi; il decadimento $\delta$
[e s$^{-1}$] e' ricavato dagli autovalori bagnati.
Lo smorzamento strutturale Rayleigh e' fissato a $\zeta=0{,}03$ per coerenza con la
Figura~5 a bassa velocita'.
Riducendo \texttt{aero\_im\_Q\_scale} sulle componenti immaginarie di $Q_{jj}$
il flutter numerico si sposta verso velocita' piu' alte
(\SI{18.1}{\knunit} con scala 0.5, \SI{31.2}{\knunit} con scala 0.4),
pur restando sotto il benchmark sperimentale.
\end{abstract}

\section{Riferimento sperimentale}

Nel documento originale il modello e' un idrofoil completamente immerso, cantilever, con
sezione NACA 16-012, massa distribuita e rigidezze pensate per rendere il sistema incline
al flutter. La Tabella~1 del paper riporta i parametri effettivi del modello; le Figure~5 e 6
forniscono rispettivamente il fattore di decadimento esponenziale e la frequenza torsionale
accoppiata in funzione della velocita'. Le tre curve teoriche del paper sono:
\begin{enumerate}
  \item coefficienti bidimensionali teorici di Theodorsen;
  \item coefficienti tridimensionali teorici di Reissner--Stevens;
  \item coefficienti tridimensionali misurati.
\end{enumerate}
La rottura sperimentale avviene a \SI{48.1}{\knunit} con frequenza accoppiata
$\approx\SI{17.5}{\cps}$.

\begin{figure}[H]
  \centering
  \begin{subfigure}{0.49\textwidth}
    \centering
    \includegraphics[width=\linewidth]{\fsipath{docs/resources/benchmarks/notes/figure_5_abramson.png}}
    \caption{Figura~5: fattore di decadimento.}
  \end{subfigure}
  \hfill
  \begin{subfigure}{0.39\textwidth}
    \centering
    \includegraphics[width=\linewidth]{\fsipath{docs/resources/benchmarks/notes/figure_6_abramson.png}}
    \caption{Figura~6: frequenza torsionale accoppiata.}
  \end{subfigure}
  \caption{Riferimenti grafici dal paper Abramson e Ransleben (1965). Nella Figura~5
  l'asse $\delta=0$ e' in alto: nel paper i valori stabili positivi compaiono verso il basso
  ($\delta\approx\SI{4}{\per\second}$ a $V=0$ sulle curve teoriche). Nei grafici numerici
  si usa la convenzione del codice (radici stabili con $\delta<0$).}
\end{figure}

\section{Verifica degli input}

La configurazione e' definita in \texttt{examples/ABRAMSON1965/config\_ABRAMSON1965.py}.
La tabella seguente confronta i principali valori fisici con quelli del documento originale.
Le conversioni sono: \SI{1}{\ftunit}=\SI{0.3048}{\meter},
\num{1} lb.in$^2$ = \SI{2.8689e-3}{\newton\meter\squared}.

\begin{table}[H]
  \centering
  \caption{Confronto tra input del paper e configurazione numerica.}
  \begin{tabular}{llll}
    \toprule
    Grandezza & Abramson 1965, valore effettivo & Config & Nota \\
    \midrule
    Semichord $b$ & \SI{0.50}{\ftunit} = \SI{0.1524}{\meter} & $c/2=\SI{0.1524}{\meter}$ & OK \\
    Chord $c$ & \SI{1.0}{\ftunit} = \SI{0.3048}{\meter} & \SI{0.3048}{\meter} & OK \\
    Semispan modello & \SI{2.5}{\ftunit} = \SI{0.762}{\meter} & \SI{0.762}{\meter} & OK \\
    Aspect ratio full-span & 5.0 & $2L/c=5.0$ & OK \\
    Massa per unita' di span & $\approx \SI{72.14}{\kilogram\per\meter}$ & \SI{72.14}{\kilogram\per\meter} & OK \\
    Mass ratio & 0.99 & $\mu/(\pi\rho b^2)\approx0.99$ & OK \\
    Asse elastico & quarto di corda & \texttt{xea\_factor}=0.25 & OK \\
    Centro di massa & $x_\alpha$ da SC & \texttt{xcm\_factor}=0.512 & OK \\
    $r_\alpha^2$ (inerzia polare) & 0.512 & 0.512 ($I_p=\mu r_\alpha^2 b^2$) & OK \\
    $EI$ & $3.40\cdot10^6$ lb.in$^2$ & \SI{9.76e3}{\newton\meter\squared} & OK \\
    $GJ$ & $0.973\cdot10^6$ lb.in$^2$ & \SI{2.79e3}{\newton\meter\squared} & OK \\
    F. flessionale a secco uncpl & \SI{10.0}{\hertz} & \SI{10.5}{\hertz} & OK \\
    F. torsionale a secco uncpl & \SI{20.5}{\hertz} & \SI{20.6}{\hertz} & OK \\
    $\zeta_\mathrm{str}$ (Rayleigh, 2 modi) & --- & 0.03 & v.\S\ref{sec:zeta} \\
    \bottomrule
  \end{tabular}
\end{table}

\section{Selezione dello smorzamento strutturale $\zeta$}
\label{sec:zeta}

Nel codice lo smorzamento strutturale e' modellato con matrice Rayleigh
$C=\alpha M+\beta K$, con coefficienti $\alpha$, $\beta$ ricavati da
\texttt{rayleigh\_target\_zetas} sui modi secchi usati nel flutter
(\texttt{rayleigh\_mode\_ids} = modi 0 e 1).
La configurazione usa \texttt{rayleigh\_target\_zetas=(0.03, 0.03)}, cioe'
$\zeta_1=\zeta_2=3\%$ sui modi di flessione e torsione a secco.

\paragraph{Allineamento con la Figura~5.}
A $V\to 0$ le curve teoriche del paper partono da
$\delta\approx\SI{4}{\per\second}$ (in convenzione codice $\delta\approx\SI{-4}{\per\second}$).
Per il modo torsionale a secco $f_2\approx\SI{20.6}{\hertz}$,
$\omega_2\approx\SI{129}{\radian\per\second}$:
\[
  |\delta| \approx \zeta\,\omega
  \quad\Rightarrow\quad
  \zeta_2 \approx \frac{|\delta|}{\omega_2}
  \approx \frac{4}{129} \approx 0.031.
\]
Il valore $\zeta=0{,}03$ e' stato scelto con questo \emph{reverse engineering} per
replicare l'ordine di grandezza del decadimento a bassa velocita' della Figura~5.

\section{Calcolo di $\delta$ dagli autovalori bagnati}
\label{sec:delta}

Per ogni velocita' $V$ il solver traccia due rami modali e risolve il problema
agli autovalori del sistema idro-aeroelastico bagnato. Per ciascun modo
con autovalore $p=\sigma+\mathrm{j}\omega$ ($\omega=|\mathrm{Im}(p)|$):
\begin{align}
  g &= \frac{\sigma}{\omega}
  \quad\text{(rapporto di smorzamento, convenzione $V$--$g$)}, \\
  \delta &= g\,\omega = \sigma
  \quad\text{(fattore di decadimento esponenziale [s$^{-1}$])}, \\
  \Lambda &= 2\pi\,g = 2\pi\,\frac{\sigma}{\omega}
  \quad\text{(decremento logaritmico per ciclo).}
\end{align}
L'implementazione e' in \texttt{vgvf\_plotting.abramson\_decay\_factor\_from\_storage};
il pannello $V$--$\delta$ di \texttt{vg\_lambda\_vf\_combined.png} esporta
\texttt{DecayFactor\_s\_inv\_Branch\_*} in \texttt{vlambda\_data.csv}.
Il flutter e' individuato al primo attraversamento di $g=0$ (equivalentemente $\delta=0$)
sul ramo che diventa instabile (tipicamente la branca~2).

\section{Metodo numerico}
\label{sec:metodo-numerico}

\subsection{Modello strutturale e aerodinamica DLM}

La run di riferimento usa un modello strutturale beam cantilever a 100 elementi, ridotto
modalmente alle prime due forme. L'inerzia polare torsionale usa $I_p=\mu\,r_\alpha^2\,b^2$
con $r_\alpha^2=0{,}512$ dalla Tabella~1. L'asse elastico e' al quarto di corda
(\texttt{xea\_factor}=0.25) e il baricentro a \texttt{xcm\_factor}=0.512.
L'aerodinamica non stazionaria e' ottenuta da matrici PanelAero/DLM precomputate su griglia
\texttt{nspan=32}, \texttt{nchord=16}, metodo \texttt{quartic}, fluido acqua e angolo
d'attacco nullo.

\subsection{Soluzione flutter con massa aggiunta strip}

Per entrambi i metodi si usa \texttt{added\_mass\_strip\_theory=True} (massa aggiunta
proiettata sui modi secchi) e, salvo diversa indicazione,
\texttt{aero\_im\_Q\_scale}=1.
Con $\zeta_\mathrm{str}=0{,}03$ le due formulazioni condividono lo stesso modello
strutturale e le stesse matrici DLM.

\paragraph{Metodo PK.}
Il metodo $P$--$K$ (\texttt{pk\_method=True}) itera sulla frequenza ridotta $k=\omega/V$
fino a convergenza del modo tracciato (MAC, predizione del passo precedente).
Fornisce le curve $V$--$g$ e $V$--$f$ usate come riferimento diretto per la validazione.

\paragraph{Metodo Roger RFA.}
La formulazione Roger (\texttt{roger\_fit=True}) approssima $Q(k)$ con stati di lag e
risolve un autovalore del sistema aumentato a ogni velocita' (senza iterazione PK su $k$).
I poli di lag sono $[0.5,\,1.0,\,1.5,\,2.0]$ ($\beta=k\,b$). Il fit Roger usa $Q(k)$ non
scalato; \texttt{aero\_im\_Q\_scale} agisce solo su $\mathrm{Im}(Q)$ in $C_\mathrm{eff}$.

\paragraph{Confronto PK+strip vs RFA+strip.}
Le Figure~\ref{fig:pk-strip} e~\ref{fig:rfa-strip} riportano le stesse grandezze
($V$--$g$, $V$--$f$; il PK produce anche il pannello $V$--$\delta$ se abilitato nel post-processing)
con \texttt{aero\_im\_Q\_scale}=1. L'accordo tra i due metodi valida l'implementazione RFA
rispetto al solver PK sullo stesso caso idro-aeroelastico.

\begin{figure}[H]
  \centering
  \includegraphics[width=\linewidth]{\fsipath{output_data/ABRAMSON1965/ABRAMSON_PK-STRIP_vgvf_combined.png}}
  \caption{Flutter PK + strip theory ($\zeta_\mathrm{str}=0{,}03$,
  \texttt{aero\_im\_Q\_scale}=1).}
  \label{fig:pk-strip}
\end{figure}

\begin{figure}[H]
  \centering
  \includegraphics[width=\linewidth]{\fsipath{output_data/ABRAMSON1965/ABRAMSON_RFA-STRIP_vgvf_combined.png}}
  \caption{Flutter Roger RFA + strip theory (stessi input del pannello PK:
  $\zeta_\mathrm{str}=0{,}03$, \texttt{aero\_im\_Q\_scale}=1).}
  \label{fig:rfa-strip}
\end{figure}

\subsection{Correzione idroelastica Capytaine (opzionale)}

Oltre alla strip theory, il codice supporta massa aggiunta e smorzamento radiativo
da Capytaine (BEM) in banda asintotica. La convenzione nel solver e':
\[
  M_\mathrm{eff} = \hat{M} + A_\infty,
  \qquad
  C_\mathrm{eff} = \hat{C} - s_Q \frac{1}{2}\rho V B_1 + B_\infty + C_\mathrm{emp},
  \qquad
  K_\mathrm{eff} = \hat{K} - \frac{1}{2}\rho V^2 B_0 ,
\]
con $s_Q=\texttt{aero\_im\_Q\_scale}$.
Per profondita' \SI{0.0508}{\meter} (2\,in sopra la lamina, come in Fig.~1 del paper),
i valori asintotici tipici sono $A_\infty=\operatorname{diag}(1.31,\,0.28)$\,kg in coordinate
modali e $B_\infty\approx 0$ nella banda di flutter.
La run RFA+strip di questo report non sostituisce la strip con Capytaine in $M_\mathrm{eff}$,
ma le curve BEM restano utili per interpretare i contributi idrodinamici.

\subsection{Smorzamento empirico}

E' disponibile uno smorzamento empirico aggiuntivo (\texttt{empirical\_fluid\_model=abramson\_delta},
\texttt{empirical\_fluid\_delta\_add\_s=5.0}\,s$^{-1}$) per modellare dissipazioni non
catturate dal DLM (viscosita', superficie libera, supporti). Nel confronto diretto con la
Figura~5 va usato con cautela perche' il paper non separa i contributi.

\section{Risultati numerici: contributi e idrodinamica modale}

Questa sezione riporta i diagrammi di supporto (contributi $M/K/C$, partecipanti DLM,
massa aggiunta) ottenuti con il solver RFA+strip a \texttt{aero\_im\_Q\_scale}=1.
Le curve di flutter con scale ridotte di $\mathrm{Im}(Q)$ sono nel
\S\ref{sec:imq-scale}; il confronto PK vs RFA e' nel \S\ref{sec:metodo-numerico}.

\begin{figure}[H]
  \centering
  \includegraphics[width=\linewidth]{\fsipath{output_data/ABRAMSON1965/stiffness_damping_contributions.png}}
  \caption{Contributi modali di massa, rigidezza e smorzamento. Convenzione:
  $M_\mathrm{eff}=M_\mathrm{struct}-M_\mathrm{aero}+M_\mathrm{hydro}$,
  $K_\mathrm{eff}=K_\mathrm{struct}-K_\mathrm{aero}$,
  $C_\mathrm{eff}=C_\mathrm{struct}-C_\mathrm{aero}+C_\mathrm{hydro}+C_\mathrm{emp}$.}
\end{figure}

\begin{figure}[H]
  \centering
  \includegraphics[width=0.95\linewidth]{\fsipath{output_data/ABRAMSON1965/ABRAMSON1965_DLM_participants.png}}
  \caption{Contributi strutturali modali $(\hat{C},\hat{K})$ e valori effettivi dopo DLM/RFA
  $(C_\mathrm{eff},K_\mathrm{eff})$.}
\end{figure}

\begin{figure}[H]
  \centering
  \includegraphics[width=0.85\linewidth]{\fsipath{output_data/ABRAMSON1965/ABRAMSON1965_Madded_terms_vs_V_capytaine_vs_strip.png}}
  \caption{Termini di massa aggiunta modale generalizzata ($m_{11}$, $m_{12}$, $m_{21}$, $m_{22}$):
  confronto tra Capytaine asintotica e strip theory proiettata sugli stessi modi.
  I valori strip sono costanti con $V$; non vanno confusi con la massa aggiunta fisica totale.}
\end{figure}

\subsection{Massa fisica totale, massa strutturale modale e massa aggiunta modale}

Nel grafico precedente compaiono termini dell'ordine di \SI{0.3}{\kilogram}--\SI{1.3}{\kilogram}.
Per interpretare correttamente questi numeri e' necessario distinguere tre livelli:
massa fisica totale dell'ala, massa strutturale modale generalizzata $\hat{M}$ e massa
aggiunta modale generalizzata $A_\infty$.

\paragraph{Massa strutturale modale generalizzata ($\hat{M}$).}
La riduzione modale produce $\hat{M}=\Phi^\top M_\mathrm{struct}\,\Phi$, con $\Phi$ costituito
dagli autovettori secchi normalizzati rispetto alla massa strutturale
($\phi_i^\top M_\mathrm{struct}\,\phi_i=1$). I termini diagonali risultano
$\hat{M}_{ii}\approx\SI{1}{\kilogram}$: sono masse di riferimento del modo $i$ nel sistema
ridotto, non la massa fisica $\mu L\approx\SI{55}{\kilogram}$ dell'idrofoil.

\paragraph{Massa aggiunta fisica totale (strip theory, moto rigido).}
\[
  m_a' = \rho\,\frac{\pi c^2}{4}\quad [\si{\kilogram\per\meter}],
  \qquad
  M_\mathrm{add,tot}^{\mathrm{strip}} \approx m_a'\,L .
\]
Con $\rho=\SI{997}{\kilogram\per\meter\cubed}$, $c=\SI{0.3048}{\meter}$, $L=\SI{0.762}{\meter}$
si ottiene $M_\mathrm{add,tot}^{\mathrm{strip}}\approx\SI{55}{\kilogram}$, comparabile a $\mu L$.

\paragraph{Massa aggiunta modale generalizzata.}
Nel solver $M_\mathrm{eff}=\hat{M}+A_\infty$ con $A_{ij}=\Phi^\top M_\mathrm{add}\,\Phi$.
Capytaine fornisce $A(\omega)$ su mesh 3D; in banda asintotica ad alta $\omega$,
per profondita' \SI{0.0508}{\meter}:
$A_{11}\approx\SI{1.31}{\kilogram}$, $A_{22}\approx\SI{0.28}{\kilogram}$,
$A_{12}=A_{21}\approx\SI{-0.41}{\kilogram}$.
Rispetto a $\hat{M}_{ii}\approx\SI{1}{\kilogram}$, sul modo~1 la sola massa aggiunta Capytaine
vale circa il 130\% di $\hat{M}_{11}$ e $M_{\mathrm{eff},11}\approx\SI{2.31}{\kilogram}$.

\begin{figure}[H]
  \centering
  \includegraphics[width=0.92\linewidth]{\fsipath{FLUID/capytaine/ABRAMSON1965/results_modal_radiation/added_mass_vs_omega_all_depths.png}}
  \caption{Massa aggiunta modale $A(\omega)$ da Capytaine per due profondita' di immersione.
  Per $\omega\gtrsim\SI{50}{\radian\per\second}$ i termini diagonali convergono ai valori
  asintotici usati nel solver idroelastico.}
  \label{fig:capytaine-added-mass-omega}
\end{figure}

\begin{figure}[H]
  \centering
  \includegraphics[width=0.92\linewidth]{\fsipath{FLUID/capytaine/ABRAMSON1965/results_modal_radiation/radiation_damping_vs_omega_all_depths.png}}
  \caption{Smorzamento radiativo modale $B(\omega)$ da Capytaine. Oltre
  $\omega\approx\SI{25}{\radian\per\second}$ il damping radiativo tende a zero; nella banda di
  flutter il solver usa $B_\infty\approx 0$.}
  \label{fig:capytaine-radiation-damping-omega}
\end{figure}

\begin{table}[H]
  \centering
  \caption{Riepilogo delle scale di massa per ABRAMSON1965 (primi due modi).}
  \label{tab:mass-scales}
  \begin{tabular}{llcc}
    \toprule
    Livello & Simbolo / origine & Modo 1 & Modo 2 \\
    \midrule
    Massa fisica totale struttura & $\mu L$ & \multicolumn{2}{c}{$\approx\SI{55}{\kilogram}$} \\
    Massa fisica totale aggiunta (strip) & $m_a'L$ & \multicolumn{2}{c}{$\approx\SI{55}{\kilogram}$} \\
    \midrule
    Massa strutturale modale & $\hat{M}_{ii}$ & \SI{1.00}{\kilogram} & \SI{1.00}{\kilogram} \\
    Massa aggiunta modale (Capytaine) & $A_{ii}$ & \SI{1.31}{\kilogram} & \SI{0.28}{\kilogram} \\
    Rapporto $A_{ii}/\hat{M}_{ii}$ (Capytaine) & -- & $\approx 130\%$ & $\approx 28\%$ \\
    Massa aggiunta modale (strip) & $\hat{M}_{\mathrm{add},ii}^{\mathrm{strip}}$ & \SI{0.26}{\kilogram} & \SI{0.45}{\kilogram} \\
    Massa effettiva modale & $M_{\mathrm{eff},ii}$ & \SI{2.31}{\kilogram} & \SI{1.28}{\kilogram} \\
    Rapporto $M_{\mathrm{eff},ii}/\hat{M}_{ii}$ & -- & $\approx 230\%$ & $\approx 128\%$ \\
    \bottomrule
  \end{tabular}
\end{table}

\section{Confronto quantitativo $\mathrm{Im}(Q)$ DLM vs coefficienti misurati (curva~3)}
\label{sec:imq-measured}

Per verificare se il flutter anticipato del codice e' coerente con un eccesso di
smorzamento circolatorio DLM rispetto ai coefficienti oscillatori \emph{misurati}
(Abramson \& Ransleben, curva~3), e' stato implementato il modulo
\texttt{COUPLING/hydroelastic\_utl/abramson\_imq\_comparison.py} con tabella
mid-span (50\,\%) da DTIC TR-4 (Tables II--X, $U=V/(b\omega)$).

\subsection{Due percorsi distinti: misura di galleria e DLM}

Abramson \textbf{non} ha misurato le matrici $Q_{jj}$ del PanelAero ne' la
$Q_\mathrm{modal}$ del solver attuale. Ha misurato \textbf{coefficienti oscillatori sezionali}
su un \textbf{modello dedicato} in galleria (Ransleben \& Abramson, JSR 1963; DTIC TR-4);
il paper di flutter del 1965 \textbf{riusa} quei coefficienti in un calcolo a strip theory
(curva~3, Figura~5). Il codice FSI calcola invece una matrice di forza aerodinamica
\textbf{generalizzata modale} $Q_\mathrm{modal}$ dal DLM su un modello FE+griglia pannelli
diverso. Il confronto nel grafico collega due mondi con una \textbf{proxy calibrata}, non
un'identita' fisica termine-per-termine.

\begin{enumerate}
  \item \textbf{Percorso Abramson (1963--65):}
  modello misura carichi $\to$ $L_{hh},L_{h\alpha},M_{hh},M_{\alpha\alpha}$ vs $U$
  $\to$ strip a 10 stazioni $\to$ curva~3 (flutter conservativo su $\delta$).
  \item \textbf{Percorso numerico:}
  DLM $Q_{jj}(k)$ $\to$ spline $Z$, $Z_{qs}$, $Z_\mathrm{force}$, $\Phi$
  $\to$ $Q_\mathrm{modal}$ (blocco $2\times2$ flessione--torsione)
  $\to$ Roger RFA $\to$ $C_\mathrm{aero}\propto\mathrm{Im}(Q)$.
\end{enumerate}

\subsection{Coefficienti misurati: come e cosa misura Abramson}

\paragraph{Esperimento di galleria (Ref.~2--3).}
Idroprofilo flessibile in \textbf{galleria chiusa}, trainato, completamente immerso,
subcavitante. Moti \textbf{imposti} armonici (non il modello flutter che poi si rompe a
\SI{48.1}{\knunit}):
\begin{itemize}
  \item \textbf{flessione} (sussulto) $h$ --- cantilever;
  \item \textbf{torsione} $\alpha$ --- intorno all'asse elastico (quarto di corda nel modello
  flutter).
\end{itemize}
Per ciascun moto e velocita' si registrano \textbf{portanza oscillante} $L$ e
\textbf{momento oscillante} $M$ (inizialmente al 31\,\% di corda, poi convertiti al
\textbf{quarto di corda} = asse elastico del modello flutter, come nel paper).

\paragraph{I quattro fattori della matrice $H(U)$.}
Sono le derivate oscillatori \textbf{2D sezionali} (notazione classica flessione--torsione).
Ogni voce e' \textbf{complessa} (modulo + fase in gradi): ampiezza e ritardo del carico
rispetto al moto armonico. Nella tabella digitata (stazione 50\,\% span) compaiono:

\begin{table}[H]
  \centering
  \caption{Coefficienti oscillatori misurati (mid-span, quarto di corda).}
  \label{tab:H-coeffs}
  \begin{tabular}{lll}
    \toprule
    Simbolo & Significato fisico & Moto che lo genera \\
    \midrule
    $L_{hh}$ & Portanza dovuta a flessione & $h$ (sussulto) \\
    $L_{h\alpha}$ & Portanza dovuta a torsione & $\alpha$ \\
    $M_{hh}$ & Momento dovuto a flessione & $h$ \\
    $M_{\alpha\alpha}$ & Momento dovuto a torsione & $\alpha$ \\
    \bottomrule
  \end{tabular}
\end{table}

Forma strip tipica (armonico, $q=\tfrac{1}{2}\rho V^2$):
\begin{align}
  L &= q\,b\left(L_{hh}\,\frac{h}{b} + L_{h\alpha}\,\alpha\right), \\
  M &= q\,b^2\left(M_{hh}\,\frac{h}{b} + M_{\alpha\alpha}\,\alpha\right).
\end{align}
Nel DTIC TR-4: \textbf{nove valori} di $U$ (0--1.67), \textbf{dieci stazioni} lungo il
semispan; per il confronto semplificato si usa la stazione \textbf{50\,\%}, coerente con
un'ala uniforme spanwise.

\paragraph{Uso nel paper 1965.}
Il \textbf{modello flutter} e' separato (trave + strip, $\mu\approx 1$).
La curva~3 della Figura~5 usa la matrice $H$ \textbf{misurata} in strip theory a 10 stazioni:
il flutter calcolato risulta \textbf{conservativo} (velocita' di instabilita' \emph{superiore}
a quella sperimentale). Il grafico $\mathrm{Im}(Q)$ di questo report \textbf{non} riproduce
quel calcolo strip: confronta un solo canale flessionale modale DLM con una proxy da $L_{hh}$.

\subsection{Matrice aerodinamica modale del codice (DLM)}

\paragraph{Catena nel solver.}
\begin{enumerate}
  \item $Q_{jj}(k)$ --- matrice aerodinamica \textbf{pannello--pannello} (DLM 3D, acqua,
  griglia \texttt{nspan=32}, \texttt{nchord=16}, NACA 16-012).
  \item Downwash modale: $(\mathrm{j}k Z + Z_{qs})\Phi$, applicazione di $Q_{jj}$, integrazione
  con aree pannelli e $Z_\mathrm{force}$.
  \item $Q_\mathrm{modal} = Z_\mathrm{force}^\top A_\mathrm{pan}\,w_\mathrm{modal}$
  (implementazione in \texttt{aero\_system.py}).
\end{enumerate}
La parte \textbf{immaginaria} di $Q$ entra nello smorzamento effettivo come
$C_\mathrm{aero}\propto (\rho V/2k)\,\mathrm{Im}(Q)$, equivalente al termine $-\tfrac{1}{2}\rho V B_1$
nel fit Roger.

\paragraph{Significato degli indici $bb$, $pp$, \ldots}
Nel blocco $2\times2$ sui primi due modi secchi:

\begin{table}[H]
  \centering
  \caption{Corrispondenza indici $Q_\mathrm{modal}$ e modi strutturali.}
  \label{tab:Q-indices}
  \begin{tabular}{llcc}
    \toprule
    Simbolo report & Modo strutturale (secco) & $f$ [Hz] & Ruolo \\
    \midrule
    $Q_{bb}$ & Modo 0, flessione & $\approx 10.5$ & Auto-influenza flessionale \\
    $Q_{pp}$ & Modo 1, torsione & $\approx 20.6$ & Auto-influenza torsionale \\
    $Q_{bp}$, $Q_{pb}$ & Accoppiamento & -- & Non nel grafico semplificato \\
    \bottomrule
  \end{tabular}
\end{table}

$Q_{bb}$ \textbf{non} e' uguale a $L_{hh}$: e' gia' una forza generalizzata modale
[N], ottenuta proiettando la risposta 3D pannello--trave su $\Phi$, non un coefficiente
sezionale adimensionale.

\subsection{Cosa confrontiamo esattamente nel grafico}

\begin{table}[H]
  \centering
  \caption{Quantita' confrontate nella Figura~\ref{fig:imq-measured}.}
  \label{tab:imq-compare}
  \begin{tabular}{p{2.2cm}p{5.2cm}p{5.5cm}}
    \toprule
    Curva & Definizione & Natura \\
    \midrule
    DLM &
    $\mathrm{Im}(Q_{bb,\mathrm{DLM}})$ da $Q_\mathrm{modal}[0,0]$ &
    Forza generalizzata; $k(V)=\omega_\mathrm{bend}(V)/V$ del ramo bagnato \\
    Misurata (cal.) &
    $\texttt{scale}_{bb}\,\mathrm{Im}(L_{hh}(U))$, $U=V/(b\omega)$ &
    Proxy: coefficiente sezionale $\times$ fattore di scala \\
    Rapporto &
    $|\mathrm{Im}(Q_\mathrm{DLM})|/|\mathrm{Im}(Q_\mathrm{meas})|$ &
    Adimensionale; evidenzia gap di ampiezza \\
    \bottomrule
  \end{tabular}
\end{table}

\textbf{Non} si confrontano: (i) $Q_{jj}$ grezze con Abramson; (ii) l'intero blocco
$2\times2$ misurato ($L_{h\alpha}$, $M_{hh}$, $M_{\alpha\alpha}$ ignorati nel grafico);
(iii) il calcolo strip completo della curva~3.

\paragraph{Convenzione di frequenza ridotta.}
PanelAero usa $k=\omega/V$ [\si{\per\meter}]. Nel paper $U=V/(b\omega)$, quindi
\begin{equation}
  k = \frac{1}{U\,b}
  \qquad\text{(non }k=1/U\text{).}
\end{equation}
Il confronto avviene allo stesso $k$ (equivalentemente allo stesso $U$) usando la
$\omega$ bagnata del ramo di flessione a ciascuna $V$.

\paragraph{Calibrazione delle unita' (proxy $Q_{bb,\mathrm{meas}}$).}
I coefficienti di galleria formano la matrice complessa
$H(U)=\{L_{hh},L_{h\alpha};M_{hh},M_{\alpha\alpha}\}$.
Non sono direttamente $Q_\mathrm{modal}$ del DLM. La stima nel grafico e' \textbf{diagonale
bend--bend}:
\[
  \mathrm{Im}(Q_{bb,\mathrm{meas}})(U)
  = \texttt{scale}_{bb}\,\mathrm{Im}(L_{hh}(U)),
  \qquad
  \texttt{scale}_{bb}
  = \frac{\bigl|\mathrm{Im}(Q_{bb,\mathrm{DLM}})\bigr|}{\bigl|\mathrm{Im}(L_{hh})\bigr|}
  \ \text{a}\ U=1.
\]
Valore numerico: $\texttt{scale}_{bb}\approx 0{,}14$.
Interpretazione: ``se il canale flessionale modale si comportasse come $L_{hh}$ misurato,
dopo un fattore unico di scala fissato a $U=1$, quanto sarebbe grande $\mathrm{Im}(Q_{bb})$?''

\paragraph{Estensione prevista.}
Roadmap: $Q_\mathrm{meas,modal}(U)\approx S\,H(U)\,S$ con matrice $S$ da calibrazione
$2\times2$ a $U=1$ (proiezione su $\Phi$, $Z_\mathrm{force}$), non solo $L_{hh}$.

\subsection{Corrispondenze fisiche e limiti del confronto}

\begin{table}[H]
  \centering
  \caption{Cosa e' analogo e cosa differisce tra misura Abramson e DLM.}
  \label{tab:analogy-limits}
  \begin{tabular}{p{3.2cm}p{4.8cm}p{4.8cm}}
    \toprule
    Aspetto & Abramson (misurato / curva~3) & Codice FSI (DLM) \\
    \midrule
    Geometria sezione & NACA 16-012 & Stessa \\
    Fluido & Acqua, immerso & Acqua \\
    Fisica & $L,M$ non stazionari da moto armonico & DLM potenziale \\
    Riduzione & $U=V/(b\omega)$ & Stesso $U$ se $k=\omega/V$ \\
    Modello & Galleria, modello carichi; strip 10 st. & Ala 3D FE + pannelli \\
    3D & Misura spanwise reale & Un numero modale diagonale \\
    Viscosita' & Reale (attutamento strumento aft) & Inviscido \\
    Accoppiamenti & Matrice $H$ completa & $Q_\mathrm{modal}$ piena; grafico solo $Q_{bb}$ \\
    Segno $\mathrm{Im}$ & Spesso $\mathrm{Im}(L_{hh})<0$ (fase) & $\mathrm{Im}(Q_{bb})>0$ nel dominio plottato \\
    \bottomrule
  \end{tabular}
\end{table}

\paragraph{Collegamento con $C_\mathrm{aero}$ e flutter.}
Nel bilancio modale, $C_\mathrm{eff}$ include $-C_\mathrm{aero}$ da $\mathrm{Im}(Q)$.
Un $|\mathrm{Im}(Q_{bb,\mathrm{DLM}})|$ circa \textbf{dieci volte} la proxy misurata a basse
$V$ (vedi sotto) implica smorzamento aerodinamico DLM dominante rispetto a
$C_\mathrm{struct}$ ($\zeta=0{,}03$) e $C_\mathrm{hydro}\approx 0$ nella run RFA+strip ---
coerente con l'attraversamento anticipato di $\delta=0$ ($V_f\approx\SI{8}{\knunit}$).

\paragraph{Risultato principale (da \texttt{ABRAMSON1965\_qmodal\_bend\_pitch.csv}).}
Per ogni velocita' dello sweep, $k(V)=\omega_\mathrm{bend}(V)/V$ e
$U(V)=V/(b\omega)$. Alla velocita' numerica di flutter ($V_f\approx\SI{4.1}{\meter\per\second}$,
$\approx\SI{8}{\knunit}$):
\[
  \frac{|\mathrm{Im}\,Q_{bb,\mathrm{DLM}}|}{|\mathrm{Im}\,Q_{bb,\mathrm{meas,cal}}|}
  \approx 7\text{--}12
  \quad (V<\SI{6}{\meter\per\second},\ \mathrm{media}\approx 10).
\]
Il DLM \textbf{sovrastima} $|\mathrm{Im}(Q)|$ di circa un ordine di grandezza rispetto ai
coefficienti misurati alla stessa $k$ --- coerente con l'instabilita' anticipata e con
\texttt{aero\_im\_Q\_scale}$\approx 0{,}1$--$0{,}15$ come ordine di grandezza di correzione.
Spesso $\mathrm{Im}(Q_\mathrm{DLM})$ e $\mathrm{Im}(Q_\mathrm{meas})$ hanno \textbf{segno opposto}
(fase del termine circolatorio diversa $\sim\SI{180}{\degree}$): il problema non e' solo l'ampiezza,
ma anche la convenzione di fase del contributo non stazionario.
Il rapporto dipende fortemente da $V$ perche' $\omega(V)$ e $k(V)$ cambiano lungo il ramo bagnato:
intorno a \SI{15}{\meter\per\second} il modulo delle due curve e' comparabile (rapporto $\approx 1$);
oltre \SI{20}{\meter\per\second} il DLM risulta \textbf{piu' debole} della proxy
(rapporto $\approx 0{,}2$--$0{,}3$). Il gap critico per $V_f$ numerico e' quindi nel regime
$V\lesssim\SI{6}{\meter\per\second}$, non su tutto lo sweep.

\begin{figure}[H]
  \centering
  \includegraphics[width=\linewidth]{\fsipath{output_data/ABRAMSON1965/ABRAMSON1965_ImQ_from_qmodal_vs_measured.png}}
  \caption{Confronto $\mathrm{Im}(Q_{bb})$ (scala lineare): DLM vs curva~3 calibrata a $U=1$.
  Destra: $|\mathrm{Im}(Q_\mathrm{DLM})|/|\mathrm{Im}(Q_\mathrm{meas})|$ vs $V$ (scala lineare);
  picco $\sim 10$ per $V<\SI{6}{\meter\per\second}$.}
  \label{fig:imq-measured}
\end{figure}

\paragraph{Script e dati.}
\begin{itemize}
  \item Tabella misurata:
    \texttt{docs/resources/benchmarks/abramson1965\_report/abramson\_measured\_coefficients\_span50.csv}
  \item Confronto completo (rigenera $Q(k)$ da Qjj):
    \texttt{examples/ABRAMSON1965/run\_imq\_vs\_measured\_comparison.py}
  \item Post-processing rapido da CSV esistente:
    \texttt{examples/ABRAMSON1965/postprocess\_imq\_from\_qmodal\_csv.py}
  \item Output:
    \texttt{output\_data/ABRAMSON1965/ABRAMSON1965\_ImQ\_from\_qmodal\_vs\_measured.csv}
\end{itemize}

\paragraph{Interpretazione e sintesi.}
\begin{itemize}
  \item \textbf{Abramson} misura in galleria le risposte aerodinamiche oscillatori sezionali
  $L_{hh},L_{h\alpha},M_{hh},M_{\alpha\alpha}$ vs $U$; il flutter del paper (curva~3) usa
  quella matrice in strip theory e risulta conservativo su $\delta(V)$.
  \item \textbf{Il codice} calcola $Q_\mathrm{modal}$ dal DLM su FE+pannelli; il grafico confronta
  solo $\mathrm{Im}(Q_{bb})$ con $\mathrm{Im}(L_{hh})$ scalato --- stesso tipo di informazione
  (parte circolatoria / smorzamento), ma quantita' fisiche diverse.
  \item Il paper classico e il DLM numerico non sono in contraddizione sulla ``non
  conservativita''' inviscida (manca smorzamento viscoso $\Rightarrow$ flutter teorico piu' tardi):
  qui il problema dominante e' una \textbf{sovrastima locale di $\mathrm{Im}(Q)$} DLM rispetto
  alla proxy da coefficienti misurati nel regime $V_f$, amplificata da $\rho$ acqua e dal fit
  Roger su tutta la griglia $k$.
\end{itemize}

\section{RFA + Capytaine (run recente)}
\label{sec:rfa-capy}

E' stata eseguita la linea \textbf{Roger RFA + strip} (\texttt{roger\_fit=True},
\texttt{added\_mass\_strip\_theory=True}, \texttt{rfa\_apply\_b2\_mass=False}) con sweep
$V\in[\SI{1}{\meter\per\second},\SI{23}{\meter\per\second}]$.
Il solver idroelastico completo \texttt{RFA\_PK\_method} (RFA + Capytaine asintotico in
$M_\mathrm{eff}$, $C_\mathrm{eff}$) e' disponibile con \texttt{RFA\_PK\_method=True}; nella
run documentata qui i contributi \texttt{M\_hydro\_Capytaine} e \texttt{C\_hydro\_Capytaine}
nei CSV risultano nulli perche' non e' attivo \texttt{RFA\_PK\_method} (solo strip in $M_\mathrm{hat}$).

\begin{table}[H]
  \centering
  \caption{Flutter numerico aggiornato (RFA+strip, $\zeta=0{,}03$, \texttt{aero\_im\_Q\_scale}=1).}
  \label{tab:flutter-recent}
  \begin{tabular}{lcc}
    \toprule
    Grandezza & Valore \\
    \midrule
    $V_F$ & \SI{4.06}{\meter\per\second} ($\approx\SI{7.9}{\knunit}$) \\
    $f_F$ (branca~2) & \SI{20.0}{\hertz} \\
    Sperimentale & \SI{48.1}{\knunit}, \SI{17.5}{\hertz} \\
    \bottomrule
  \end{tabular}
\end{table}

Il confronto con Capytaine asintotico ($A_\infty$ modo~1 $\approx\SI{1.31}{\kilogram}$ vs
strip $\approx\SI{0.29}{\kilogram}$) resta rilevante per la massa, ma nella banda di flutter
$B_\infty\approx 0$: il gap dominante per $V_f$ resta $\mathrm{Im}(Q)$ DLM
(\S\ref{sec:imq-measured}), non il damping radiativo Capytaine.

\section{Scala dello smorzamento idrodinamico $\mathrm{Im}(Q_{jj})$}
\label{sec:imq-scale}

Le componenti immaginarie di $Q_{jj}$ alimentano il termine di smorzamento aerodinamico
$C_\mathrm{aero}\propto \mathrm{Im}(Q)$. Riducendo \texttt{aero\_im\_Q\_scale} si attenua
lo smorzamento idrodinamico DLM senza alterare il fit Roger su $\mathrm{Re}(Q)$.
Le figure seguenti usano sempre RFA+strip e $\zeta_\mathrm{str}=0{,}03$.

\begin{figure}[H]
  \centering
  \includegraphics[width=\linewidth]{\fsipath{output_data/ABRAMSON1965/vdelta_vf_Imq_scaled_05_RFA.png}}
  \caption{RFA+strip, \texttt{aero\_im\_Q\_scale = 0.5}: $V_f\approx\SI{18.1}{\knunit}$
  sulla branca~2.}
  \label{fig:imq-05}
\end{figure}

\begin{figure}[H]
  \centering
  \includegraphics[width=\linewidth]{\fsipath{output_data/ABRAMSON1965/vdelta_vf_Imq_scaled_04_RFA.png}}
  \caption{RFA+strip, \texttt{aero\_im\_Q\_scale = 0.4}: $V_f\approx\SI{31.2}{\knunit}$
  sulla branca~2.}
  \label{fig:imq-04}
\end{figure}

\paragraph{Discussione.}
Con \texttt{aero\_im\_Q\_scale}$=1$ la branca~2 attraversa l'instabilita' intorno a
$V_f\approx\SI{7.9}{\knunit}$ (Tab.~\ref{tab:flutter-recent}).
Con fattore \textbf{0.5}, $\delta(V)$ cresce piu' lentamente e $V_f\approx\SI{18.1}{\knunit}$
($f_F\approx\SI{20.7}{\hertz}$).
Con fattore \textbf{0.4}, $V_f\approx\SI{31.2}{\knunit}$ ($f_F\approx\SI{17.5}{\hertz}$),
vicino alla frequenza sperimentale ma ancora sotto i \SI{48.1}{\knunit} di rottura.
Il ramo~1 resta stabilmente smorzato su tutto il dominio plottato.
La sensibilita' conferma che il bilancio tra smorzamento strutturale ($\zeta=0{,}03$) e
smorzamento idrodinamico DLM governa la posizione del flutter.

\section{Confronto con Figure~5 e~6 ridigitalizzate}

I punti sperimentali e la curva teorica~3 (coefficienti 3D misurati) sono stati
ridigitalizzati dai PNG in \texttt{docs/resources/benchmarks/notes/} e salvati in
\texttt{abramson\_fig5\_digitized.csv}, \texttt{abramson\_fig6\_digitized.csv} e
\texttt{abramson\_rfa\_strip\_selected.csv} (curve RFA+strip, \texttt{aero\_im\_Q\_scale}=0.5).
I valori sono approssimati (lettura grafica); servono al trend qualitativo.

\begin{figure}[H]
  \centering
  \begin{tikzpicture}
    \begin{axis}[
      width=0.95\linewidth,
      height=0.52\linewidth,
      xlabel={Velocita' $V$ [kn]},
      ylabel={$\delta$ [s$^{-1}$]},
      xmin=0, xmax=55,
      ymin=-15, ymax=5,
      unbounded coords=jump,
      legend style={font=\small, at={(0.02,0.02)}, anchor=south west}
    ]
      \addplot+[only marks, mark=o, black] table[x=Speed_knots,y=Exp_decay_s_inv,col sep=comma]
        {\fsipath{docs/resources/benchmarks/abramson1965_report/abramson_fig5_digitized.csv}};
      \addlegendentry{Abramson, dati sperimentali}
      \addplot+[only marks, mark=square, gray] table[x=Speed_knots,y=Curve3_calc_decay_s_inv,col sep=comma]
        {\fsipath{docs/resources/benchmarks/abramson1965_report/abramson_fig5_digitized.csv}};
      \addlegendentry{Abramson, curva 3 (calc.)}
      \addplot+[very thick, red] table[x=Speed_knots,y=Decay_s_inv_Branch_2,col sep=comma]
        {\fsipath{docs/resources/benchmarks/abramson1965_report/abramson_rfa_strip_selected.csv}};
      \addlegendentry{RFA+strip, branca 2}
      \addplot+[thick, red!45, dashed] table[x=Speed_knots,y=Decay_s_inv_Branch_1,col sep=comma]
        {\fsipath{docs/resources/benchmarks/abramson1965_report/abramson_rfa_strip_selected.csv}};
      \addlegendentry{RFA+strip, branca 1}
      \addplot[densely dashed, black] coordinates {(48.1,-15) (48.1,5)};
      \addplot[densely dashed, red] coordinates {(18.1,-15) (18.1,5)};
    \end{axis}
  \end{tikzpicture}
  \caption{Confronto $\delta(V)$: dati e curva~3 ridigitalizzati vs RFA+strip
  ($\zeta=0{,}03$, \texttt{aero\_im\_Q\_scale}=0.5).}
\end{figure}

\begin{figure}[H]
  \centering
  \begin{tikzpicture}
    \begin{axis}[
      width=0.95\linewidth,
      height=0.52\linewidth,
      xlabel={Velocita' $V$ [kn]},
      ylabel={Frequenza [Hz]},
      xmin=0, xmax=55,
      ymin=14, ymax=25,
      unbounded coords=jump,
      legend style={font=\small, at={(0.98,0.02)}, anchor=south east}
    ]
      \addplot+[only marks, mark=o, black] table[x=Speed_knots,y=Exp_frequency_cps,col sep=comma]
        {\fsipath{docs/resources/benchmarks/abramson1965_report/abramson_fig6_digitized.csv}};
      \addlegendentry{Abramson, dati sperimentali}
      \addplot+[only marks, mark=square, gray] table[x=Speed_knots,y=Curve3_calc_frequency_cps,col sep=comma]
        {\fsipath{docs/resources/benchmarks/abramson1965_report/abramson_fig6_digitized.csv}};
      \addlegendentry{Abramson, curva 3 (calc.)}
      \addplot+[very thick, red] table[x=Speed_knots,y=Frequency_Hz_Branch_2,col sep=comma]
        {\fsipath{docs/resources/benchmarks/abramson1965_report/abramson_rfa_strip_selected.csv}};
      \addlegendentry{RFA+strip, branca 2}
      \addplot+[thick, red!45, dashed] table[x=Speed_knots,y=Frequency_Hz_Branch_1,col sep=comma]
        {\fsipath{docs/resources/benchmarks/abramson1965_report/abramson_rfa_strip_selected.csv}};
      \addlegendentry{RFA+strip, branca 1}
      \addplot[densely dashed, black] coordinates {(48.1,14) (48.1,25)};
    \end{axis}
  \end{tikzpicture}
  \caption{Confronto frequenza accoppiata.}
\end{figure}

\begin{table}[H]
  \centering
  \caption{Riepilogo flutter (branca instabile, RFA + strip, $\zeta=0{,}03$).}
  \begin{tabular}{lccc}
    \toprule
    Caso & \texttt{aero\_im\_Q\_scale} & $V_F$ [kn] & $f_F$ [Hz] \\
    \midrule
    Sperimentale & --- & 48.1 & 17.5 \\
    RFA + strip & 1.0 & $\sim$8 & $\sim$20 \\
    RFA + strip & 0.5 & 18.1 & $\approx$20.7 \\
    RFA + strip & 0.4 & 31.2 & $\approx$17.5 \\
    \bottomrule
  \end{tabular}
\end{table}

\section{Conclusioni}

La pipeline produce modi secchi, matrici DLM, sweep Roger RFA, contributi $M/K/C$ e grafici
$V$--$\delta$--$g$--$f$ in unita' Abramson.
Le frequenze a secco (\SI{10.5}{\hertz} flessione, \SI{20.6}{\hertz} torsione) sono allineate
alla Tabella~1; le frequenze bagnate seguono il trend sperimentale, ma il flutter numerico
resta \emph{anticipato} rispetto ai \SI{48.1}{\knunit} misurati: il confronto
$\mathrm{Im}(Q)$ (\S\ref{sec:imq-measured}) indica una sovrastima DLM di circa
$\times 7$--$\times 12$ su $|\mathrm{Im}(Q_{bb})|$ alla stessa $k$ dei coefficienti misurati
nel regime $V_f$ numerico. Calibrazioni possibili: \texttt{aero\_im\_Q\_scale},
coefficienti misurati al posto del DLM, o Capytaine + strip coerente in $M_\mathrm{eff}$.
Come benchmark ABRAMSON1965, la configurazione \textbf{RFA + strip} con
$\zeta_\mathrm{str}=0{,}03$ e $\delta$ dagli autovalori bagnati e' la base di confronto
con le Figure~5 e~6 del paper.

\end{document}
